\documentclass[10pt]{article}
\usepackage[T1]{fontenc}
\usepackage[utf8]{inputenc}
\usepackage[french]{babel}
\usepackage{amsmath,amssymb,mathrsfs,mathpazo}
\usepackage{geometry}
\geometry{a4paper,margin=2cm}
\usepackage{graphicx,array,tabularx,tabularray,colortbl,multirow,multicol,booktabs,enumitem,pifont,ulem,fancybox,diagbox,cancel}
\usepackage{pstricks,pstricks-add,pst-plot,pst-tree,pst-node,pst-text,pst-func,pst-math,pst-bspline,pst-3dplot,pst-eucl,pst-solides3d}
\let\pscalculate\undefined
\usepackage{pst-bezier}
\usepackage{tikz,pgfplots,tkz-tab}
\usetikzlibrary{arrows,arrows.meta,patterns,calc,positioning,trees,decorations.markings,shapes.geometric,shapes.misc,shapes.arrows,intersections,decorations.pathmorphing,decorations.pathreplacing,backgrounds,shadows,babel}
\pgfplotsset{compat=1.18}
\usepgfplotslibrary{fillbetween}
\usepackage[np]{numprint}
\usepackage{listings,makecell,pas-tableur,twoopt}
\let\barre\undefined
\usepackage[pstricks]{tablvar}
\definecolor{darkgreen}{RGB}{0,100,0}
\providecommand{\red}{\color{red}}
\providecommand{\blue}{\color{blue}}
\providecommand{\green}{\color{green}}
\providecommand{\black}{\color{black}}
\providecommand{\white}{\color{white}}
\providecommand{\gray}{\color{gray}}
\providecommand{\cyan}{\color{cyan}}
\providecommand{\magenta}{\color{magenta}}
\providecommand{\yellow}{\color{yellow}}
\providecommand{\violet}{\color{violet}}
\providecommand{\orange}{\color{orange}}
\providecommand{\euro}{\texteuro}
\providecommand{\EUR}{\texteuro}
\providecommand{\ieme}{\textsuperscript{e}}
\providecommand{\degres}{\ensuremath{^\circ}}
\providecommand{\degree}{\ensuremath{^\circ}}
\providecommand{\vv}[1]{\overrightarrow{#1}}
\providecommand{\Coord}[1]{#1}
\providecommand{\decosix}{}
\providecommand{\starredbullet}{\ensuremath{\bullet}}
\pagestyle{empty}
\setlength{\parindent}{0pt}
\usepackage{adjustbox}\DeclareUnicodeCharacter{2236}{:}\DeclareUnicodeCharacter{2190}{\ensuremath{\leftarrow}}
\providecommand{\R}{\mathbb{R}}
\providecommand{\N}{\mathbb{N}}
\providecommand{\Z}{\mathbb{Z}}
\providecommand{\Q}{\mathbb{Q}}
\providecommand{\D}{\mathbb{D}}
\providecommand{\C}{\mathbb{C}}
\providecommand{\e}{{\mathrm{e}}}
\providecommand{\vect}[1]{\overrightarrow{#1}}
\providecommand{\vectt}[1]{\overrightarrow{\mathrm{#1}}}
\providecommand{\vvt}[1]{\overrightarrow{\mathrm{#1}}}
\providecommand{\np}[2][]{{#2~#1}}
\providecommand{\vv}[1]{{\overrightarrow{#1}}}
\renewcommand{\d}{{\mathrm{d}}}
\providecommand{\degres}{{^\circ}}
\providecommand{\degre}{{^\circ}}
\providecommand{\blue}{\color{blue}}
\providecommand{\red}{\color{red}}
\providecommand{\enskip}{\quad}
\providecommand{\llbracket}{\lbrack\!\lbrack}
\providecommand{\rrbracket}{\rbrack\!\rbrack}
\providecommand{\euro}{{\text{€}}}
\providecommand{\no}{{\text{n°}}}
\providecommand{\interval}[2]{{[#1;#2]}}
\providecommand{\intervFO}[2]{{[#1;#2[}}
\providecommand{\dotfill}{\ldots}

\providecommand{\ds}{\displaystyle}
\providecommand{\textding}[1]{#1}
\providecommand{\Oij}{\(\mathrm{O};\vec{\imath},\vec{\jmath}\)}
\providecommand{\Oijk}{\(\mathrm{O};\vec{\imath},\vec{\jmath},\vec{k}\)}
\providecommand{\Ouv}{\(\mathrm{O};\vec{u},\vec{v}\)}
\providecommand{\intervFF}[2]{\ensuremath{[#1;#2]}}
\renewcommand{\euro}{\ifmmode\text{\texteuro}\else\texteuro\fi}
\DeclareUnicodeCharacter{00B0}{\ifmmode^{\circ}\else\textdegree\fi}
\DeclareUnicodeCharacter{20AC}{\ifmmode\text{\texteuro}\else\texteuro\fi}
% Énoncés : Éducation nationale. Transcription et corrigés : APMEP, Denis Vergès et contributeurs. Découpage et classement : CoopMaths, Sébastien Lozano, Éric Elter et contributeurs. Adaptation de la présentation : Maths Explorer.
% CC BY-SA : https://coopmaths.fr/www/about/
% https://coopmaths.fr/alea/static/bac/2023/tex/bac_2023_09_sujet2_ameriquesud_3_cor.tex
\begin{document}
\noindent\textit{APMEP / CoopMaths. CC BY-SA. Présentation adaptée par Maths Explorer.}\par\medskip
\medskip

Soit la suite $\left(u_n\right)$ définie par $u_0 = 0$ et, pour tout $n \in \N$, 
$u_{n+1} = 5u_n - 8n + 6$.



\begin{enumerate}
\item
\begin{list}{\textbullet}{}
\item Pour $n=0$ on a:
$u_1=u_{0+1}=5u_0 -8\times 0 +6=0-0+6=6$
\item Pour $n=1$ on a: $u_2=u_{1+1}=5u_1 -8\times 1 +6= 5\times 6 -8+6=28$
\end{list}

\item Soit $n$ un entier naturel.
On complète la fonction {\ttfamily suite\_u} d'argument {\ttfamily n}  ci-dessous, écrite en langage Python, afin qu'elle retourne la valeur de $u_n$.

\begin{center}
\begin{tabular}{>{\ttfamily}l}
def suite\_u(n) : \\
\qquad u = \blue 0\\
\qquad for i in range(1,n+1) :\\
\qquad \vline \qquad u = \blue 5*u - 8*(i-1) + 6\\
\qquad return u
\end{tabular}
\end{center}

\emph{Rem}: il aurait sans doute été plus simple de donner comme algorithme:

\begin{center}
\begin{tabular}{>{\ttfamily}l}
def suite\_u(n) : \\
\qquad u = \blue 0\\
\qquad for i in \red range(0,n) :\\
\qquad \vline \qquad u = \blue 5*u - 8*i + 6\\
\qquad return u
\end{tabular}
\end{center}

\item
	\begin{enumerate}
		\item Soit $\mathcal{P}_n$ la propriété:  $u_n \geqslant 2n$.

\begin{list}{\textbullet}{}
\item \textbf{Initialisation}

Pour $n=0$, on a: $u_0=0$ et $2n=0$; donc $\mathcal{P}_0$ est vraie.
\item \textbf{Hérédité}

On suppose que $u_n \geqslant 2n$ pour $n\geqslant 0$; c'est l'hypothèse de récurrence.

$u_{n+1}=5u_n -8n+6 \geqslant 5\times (2n) -8n+6 = 2n+6 = 2(n+1)+4 \geqslant 2(n+1)$

Donc $\mathcal{P}_{n+1}$ est vérifiée.
\item \textbf{Conclusion} 

La propriété est vraie au rang 0, et elle est héréditaire pour tout $n\geqslant 0$, donc, d'après le principe de récurrence, elle est vraie pour tout $n$.
\end{list}		

On a donc démontré que, pour tout entier naturel $n$, on a: $u_n\geqslant 2n$.		
		
		\item
$\ds\lim _{n\to +\infty} 2n=+\infty$ et $u_n\geqslant 2n$ donc $\ds\lim _{n\to +\infty} u_n = +\infty$.
		
		\item Soit $p \in \N^*$.
		
$\ds\lim _{n\to +\infty} u_n = +\infty$ donc, d'après la définition de la limite vers $+\infty$ d'une suite, pour tout réel $A>0$, il existe un rang $n_0$ à partir duquel tous les termes de la suite sont supérieurs à $A$; donc il existe un rang $n_0$ tel que si $n\geqslant n_0$, alors $u_{n_0}\geqslant 10^p$.
\end{enumerate}

\item 
Pour tout $n$, on a:
$u_{n+1}-u_n=5u_n-8n+6 -u_n = 4u_n - 8n +6$;
or $u_n\geqslant 2n$ donc $4u_n-8n\geqslant 0$ et donc $u_{n+1}-u_n\geqslant 6$.

On en déduit que la suite $(u_n)$ est croissante.

\item On considère la suite $\left(v_n\right)$, définie pour tout $n \in  \N$, par $v_n = u_n - 2n + 1$.
	\begin{enumerate}
		\item En dessous de la fonction {\ttfamily suite\_u} précédente, on a écrit la fonction {\ttfamily suite\_v} ci-dessous :
\begin{center}
\begin{tabular}{>{\ttfamily}l}
def suite\_v(n):\\
\qquad L = []\\
\qquad  for i in range(n+1) :\\
\qquad \vline \qquad L.append(suite\_u(i) - 2*i + 1)\\
\qquad return L
\end{tabular}
\end{center}



Lorsqu'on saisit {\ttfamily suite\_v(5)} dans la console, on obtient l'affichage suivant :

\begin{center}
\begin{tabular}{>{\ttfamily} l}
$>>>$ suite\_v(5)\\
$[1,~5,~25,~125,~625,~3125]$\\
\end{tabular}
\end{center}



On peut conjecturer que, pour tout $n$, on a: $v_{n+1}=5v_n$.



$\begin{gathered}v_{n+1}=u_{n+1}-2(n+1)+1 = 5u_n -8n +6 -2n-2+1=5u_n-10n+5= 5(u_n-2n+1)\\
\phantom{v_{n+1}}
=5v_n\end{gathered}$

		\item 
On déduit de la question précédente que  la suite $(v_n)$ est géométrique de raison $q=5$, et de premier terme $v_0=u_0-2\times 0+1=1$. 

On a donc, pour tout $n$: $v_n = v_0\times q^n = 5^n$.

Pour tout $n$, $v_n=u_n-2n+1$ donc $u_n=v_n+2n-1$ et donc $u_n= 5^n +2n-1$.
	\end{enumerate}
\end{enumerate}

\bigskip

\end{document}
